\chapter{Кубический следовой связный оператор трёх копий}
\label{ch:v8-baryon-cubic-trace-connected-operator-gate}

Предыдущая глава доказала, что двухточечный процесс не определяет третью
кумулянту. Однако полный шумовой кадр уже несёт следовую метрику. Проверим,
порождает ли она канонический связный трёхкопийный оператор до введения
координатной потенциальной модели.

\section{Каноническое удаление центральной части}

Для полного самосопряжённого кадра $\{F_a\}_{a=1}^{42}\subset M_{21}$
введём бесследовую часть
\begin{equation}
 \widehat F_a=F_a-\frac{\Tr F_a}{21}I_{21}.
 \label{eq:v8-baryon-centered-noise-frame}
\end{equation}
Только один исходный оператор имеет ненулевой след, равный $-4$.
Центрирование не меняет диссипативный двойной коммутатор:
\begin{equation}
 [\widehat F_a,[\widehat F_a,X]]=[F_a,[F_a,X]].
 \label{eq:v8-baryon-centering-preserves-generator}
\end{equation}
Следовательно, бесследовый кадр задаёт ту же квантово-марковскую
полугруппу. Его грамова матрица
\begin{equation}
 \widehat K_{ab}=\Tr(\widehat F_a\widehat F_b)
 \label{eq:v8-baryon-centered-trace-metric}
\end{equation}
имеет точный ранг $42$ и допускает обратную метрику $\widehat K^{ab}$.

\section{Кубический следовой тензор}

Определим полностью симметричный вещественный тензор
\begin{equation}
 d_{abc}=\frac12\Tr\!\left(
 \widehat F_a\{\widehat F_b,\widehat F_c\}\right).
 \label{eq:v8-baryon-cubic-trace-tensor}
\end{equation}
Цикличность следа доказывает его инвариантность при сопряжённом действии
калибровочной группы. Точный аудит на всём кадре даёт $168$ ненулевых
неупорядоченных компонент с разложением
\begin{equation}
 N_{TTT}=0,\qquad N_{TTG}=140,\qquad
 N_{TGG}=0,\qquad N_{GGG}=28,
 \label{eq:v8-baryon-cubic-trace-support}
\end{equation}
где $T$ обозначает направление переноса, а $G$ --- калибровочное
направление. Таким образом, чистого кубического переноса нет, но существует
ненулевой канал «два переноса плюс калибровочное направление».

\section{Связный оператор}

Поднимем индексы тензора обратной следовой метрикой и положим
\begin{equation}
 W_3=d^{abc}\widehat F_a\otimes\widehat F_b\otimes\widehat F_c,
 \qquad
 d^{abc}=\widehat K^{aa'}\widehat K^{bb'}\widehat K^{cc'}d_{a'b'c'}.
 \label{eq:v8-baryon-connected-cubic-operator}
\end{equation}
Линейная независимость кадра и ненулевость $d$ дают $W_3\ne0$.
Вещественность и симметрия $d$ делают оператор самосопряжённым и
перестановочно-инвариантным. Поскольку все $\widehat F_a$ бесследовы,
\begin{equation}
 \Tr_1W_3=\Tr_2W_3=\Tr_3W_3=0.
 \label{eq:v8-baryon-connected-cubic-partial-traces}
\end{equation}
Значит, $W_3$ является настоящим связным трёхкопийным оператором: он
невидим во всех одно- и двухчастичных ограничениях.

\begin{theorem}[Допуск канонического связного носителя]
Полный 42-мерный шумовой кадр и его конечный след канонически определяют
ненулевой, самосопряжённый, калибровочно- и перестановочно-инвариантный
трёхкопийный оператор с нулевыми частичными следами. Для его определения не
нужны наблюдаемые массы, пространственная волновая функция или ручной выбор
ковариационного параметра.
\end{theorem}

\section{Граница физического статуса}

Текущий родитель содержит квадратичную форму Дирихле и слабый предел второго
порядка. Поэтому существование $W_3$ в алгебре допустимых операторов ещё не
означает, что действие содержит член
\begin{equation}
 H_{(3)}=\lambda_3W_3.
 \label{eq:v8-baryon-cubic-parent-term}
\end{equation}
Ни коэффициент $\lambda_3$, ни знак, ни пространственное продолжение, ни
неприводимое ядро Фаддеева текущим расчётом не выведены.

\begin{nogocriterion}[Запрет преждевременного физического чтения]
Кубический следовой тензор нельзя объявлять барионным гамильтонианом только
по его каноничности. Следующий гейт обязан проверить происхождение
$\lambda_3W_3$ из того же родительского действия. Кроме того, отсутствие
компоненты $TTT$ запрещает толковать оператор как три независимых переносных
ребра без участия калибровочного канала.
\end{nogocriterion}

\begin{protocol}
\begin{enumerate}
 \item поле вычислений: \textbf{$\mathbb Q$};
 \item числа с плавающей точкой: \textbf{отсутствуют};
 \item ранг центрированного кадра: \textbf{$42$};
 \item диссипативный генератор после центрирования: \textbf{тот же};
 \item ненулевые компоненты $d_{abc}$: \textbf{$168$};
 \item опора: \textbf{$140\,TTG+28\,GGG$};
 \item чистая $TTT$-компонента: \textbf{отсутствует};
 \item частичные следы $W_3$: \textbf{нулевые};
 \item канонический связный оператор: \textbf{допущен};
 \item коэффициент родительского действия: \textbf{не выведен};
 \item следующий гейт:
       \textbf{происхождение кубического члена либо строгий запрет}.
\end{enumerate}
\end{protocol}

Машинный сертификат сохранён в
\path{s2t_v8_baryon_cubic_trace_connected_operator_gate_results.json}.